\documentclass{article}
\usepackage[utf8]{inputenc}
\usepackage{color}
%\title{\textbf{Speaker Recognition using Machine Learning: A Study on Efficacy of Gaussian Mixture Models}}
\title{\textbf{Efficacy of a Gaussian Mixture Model based Speaker Recognition in Machine Learning Framework}}
\author{\textbf{Adhiraj Banerjee}}
%\centering{\Large{4 week report}}
%(SRFP Application No. : ENGS589)

%\date{18 July 2021}

\begin{document}

\maketitle

\section{\textbf{{Introduction}}}
%The topic for my Internship which was assigned to me by my guide is, \textit{Speaker Recognition or Identification using Machine Learning}. Here, the problem I have to solve is that I have to create a trained model which will be able to recognize the speaker from its voice sample provided to it as its input. 
Speaker Recognition (SR) has found numerous applications such as voice biometric, forensics, traitor finding \textcolor{red}{(Adi, Give a few more here)}. Conceptually, SR determines the speaker among a set of registered speakers. The present work studies available methods for speaker recognition. Further, in view of the applicability of \textbf{Machine Learning (ML)} to a wide variety of applications, in the present work emhasis is laid on the popular Gaussian Mixture Model (GMM) based ML approach as a potential solution for SR. An accuracy of $97\%$ (\textcolor{red}{provide proper number here}) is observed for GMM-ML based SR when tested over $4$ sets of speakers.

\section{\textbf{{``Machine Learning'' for Speaker Recognition}}}
% \textit{Machine Learning} is a method which helps systems to learn from the training data so that they can predict the class of the test data given as input, without being explicitly programmed for this purpose. This same method would be applied in the Speaker Recognition System. Using Machine Learning, a model for Speaker Recognition will be trained in such a way that it will be able to identify unique patterns and features from different voice samples and distinguish a voice sample from the rest. This method can be applied by, firstly \textbf{gathering different people's speech samples}, \textbf{extract features from the voice sample} suitable for the classifier, which is \textbf{trained for building a trained model} and finally,\textbf{performing classification for recognition of each Speaker} from its corresponding input voice sample.
 \textit{Machine Learning} (ML) is a mathematical framework which helps systems to learn from trained model and predict the class into which the test data may be classified into. 
%Without being explicitly programmed for this purpose. 
The present work studies Speaker Recognition from a ML perspective. Going by the principles of ML, a model for Speaker Recognition is trained in such a way that it will be able to identify unique patterns and features from different voice samples and distinguish a voice sample from the rest. This method can be applied by, firstly {gathering different people's speech samples}, {extract features from the voice sample} suitable for the classifier, which is {trained for building a trained model} and finally,{performing classification for recognition of each Speaker} from its corresponding input voice sample.

\subsection{Gaussian Mixture Model based SR}
Write a few lines on GMM and how we obtain MFCC for ML based SR.

\section{\textbf{{Results and Inference}}}
\textcolor{red}{Write a few lines on the results you have observed from your python codes. Provide quantitative results like $97\%$ etc.}

\section{\textbf{{Conclusion on Work Done}}}
%I would like to conclude this report by stating that, the 
Speech Recognition is a \textit{text-dependent} method, where it highly depends on the language and corpus, whereas, Speaker Recognition mainly focusses on raw audio percepts (\textcolor{red}{check this terminology}) and its related data by which it identifies the uniqueness among different speakers. \textit{ML} approach is a practical solution which renders \textit{Speaker} Recognition advantageous to \textit{Speech} Recognition. 
%Hence, \textbf{Machine Learning} helps us to achieve the goal of Speaker Recognition and makes us understand about differences between \textit{Speaker Recognition} \& \textit{Speech Recognition}.
\textcolor{red}{Write a few lines on the results you have observed from your python codes. Conclusion means, what you want to say finally. Keep it simple but technically appropriate to the context.}

\end{document}